\runninghead{The Hour to Rise}
\properlane{hour-to-rise}{The Hour to Rise: The Sleeper Is Told to Wake, and the One Who Comes Is Asked to Set Him Free}

St Paul tells the Romans \latin{hora est iam nos de somno súrgere}, ``it is
now the hour for us to rise from sleep,'' and gives the reason: \latin{Nunc
enim própior est nostra salus, quam cum credídimus}, ``For now our salvation
is nearer than when we believed'' (Rom 13:11). Two questions follow from the
verse. What is the salvation that is nearer, and nearer in what way? And how
does a sleeper obey a command to rise? The Epistle addresses the command to
the hearers. Before it, the Collect has asked Christ to stir up his power and
come, so that \latin{te protegénte}, with him protecting, we may be snatched
from the dangers of our sins, and \latin{te liberánte}, with him setting
free, we may be saved. This interpretation follows St John Chrysostom's
homily on the passage and St Thomas Aquinas's lecture on it.

\subsection{Chrysostom: runners near the goal}

Chrysostom reads Paul's ``hour'' as the hour of the end. ``It is, that near is
the Resurrection, near the awful Judgment, and the day that burneth as a
furnace, near.'' Every day brings it closer: ``as the time advances \dots\ the
season of our present life is wasting away, and that of the life to come
waxes nearer'' (\work{Homilies on Romans} 24, on v.~11). He notices that Paul
urges the Romans ``not upon alarming grounds \dots\ but upon kindly ones,'' and
he sees a particular danger. Converts are most eager at the start, ``but \dots\
as the time went on, the whole of their earnestness would wither down to
nothing.'' Paul's argument reverses that decline. ``The nearer the prize is,
the more wide awake ought they to be for the contest, since even the racers do
this, when they are upon the end of the course \dots\ then they rouse
themselves up the more.''

The night far spent becomes, in Chrysostom's hands, a household before dawn.
``When we see the night pressing on towards the morning, and hear the swallow
twittering, we each of us awake our neighbor, although it be night still. But
so soon as it is actually departing, we hasten one another, and say It is day
now! and we all set about the works of the day, dressing, and leaving our
dreams.'' So Christians are to ``put off imaginings \dots\ get clear of the
dreams of this life present'' and be clothed in virtue (on v.~12). He tells his
hearers not to fear the armour of light. ``It is not in fact war, but a solemn dance
and feast-day,'' and the man armed with the arms of light goes out ``as the
bridegroom goes forth with joyous looks from his chamber \dots\ For he is at
once soldier and bridegroom.''

On v.~13 he notices that Paul says ``let us walk,'' not ``walk ye,'' ``so
making the exhortation free from anything grating, and the reproof gentle,''
and that Paul forbids ``not the enjoying of wine, but doing it to excess.'' On
v.~14 he notices a change of word. When Paul spoke of vice he spoke of
``works''; when he speaks of virtue he speaks of ``arms''; and at last ``he gives
us the Lord Himself for a garment, the King Himself.'' To put on Christ is to
be ``girt about with Him upon every side,'' since ``He would have our soul to
be a dwelling for Himself, and Himself to be laid round about us as a
garment, that He may be unto us all things both from within and from
without'' (on v.~14).

\subsection{Aquinas: three ways salvation is near}

Aquinas asks what the sleep is. It is not the sleep of nature, nor the rest of
contemplation, but the sleep of sin, as in ``Rise, thou that sleepest, and
arise from the dead'' (Eph 5:14), or of negligence, as in Proverbs' ``How
long wilt thou sleep, O sluggard?'' (Prov 6:9). One rises from the first by
penance and from the second by care for good works (\work{Super Epistolam ad
Romanos} 13, lect.~3).

He then explains the nearness. Paul's \latin{salus} is eternal life, and it
is near in two ways: in time, as the just come closer to the end of their
lives, and by disposition, as charity grows and makes a person readier for
the wedding feast (Mt 25:10). That is the Apostle's own meaning. Aquinas adds
a third sense and says whose it is: \latin{secundum quod Ecclesia haec verba
adventus tempore legit}, ``as the Church reads these words in the time of
Advent,'' they speak of the salvation of Christ's first coming, the Apostle
speaking \latin{ex persona omnium fidelium qui fuerunt ab initio mundi}, ``in
the person of all the faithful who have been since the beginning of the
world.'' He cites Isaias: ``my salvation is near to come'' (Isa 56:1).

He reads v.~12 in the same order. There are three nights to which
\latin{nox præcéssit} can apply: the ignorance of this life, the night of sin,
and the time before the Incarnation. The day at hand is, by the Apostle's
intention, the day of glory; it can also be the day of grace, which has come
and still draws near to each \latin{per fidem et devotionem}, ``by faith and
devotion,'' as Paul writes elsewhere, ``The Lord is nigh'' (Phil 4:5).

\subsection{Who sets the sleeper free}

Paul commands the Romans to rise. St Augustine, in the \work{Confessions},
describes what it is like to hear such a command and not obey it. Before his
conversion he was ``like unto the efforts of those desiring to awake, who,
still overpowered with a heavy drowsiness, are again steeped therein.'' A man,
he says, puts off shaking off sleep and enjoys it the more though it
displeases him, \latin{quamvis surgendi tempus advenerit}, ``although the
time for rising has come.'' He had no answer to the call ``Awake, thou that
sleepest,'' only ``the drawling and drowsy words: `Presently, lo, presently;'
`Leave me a little while.' But `presently, presently,' had no present.'' His
own will could not free him: ``who shall deliver me from the body of this
death'' but ``Thy grace only, through Jesus Christ our Lord?''
(\work{Confessions} VIII.5.12). When he was freed, it was by these verses of
Romans (VIII.12.29).

His story shows why a command to wake is not enough by itself: a man may know
it is the hour and still say ``presently.'' His question is St Paul's (Rom
7:24), made his own: \latin{miserum ergo me quis liberaret}, ``who then would
set me free, wretched as I am?'' The Collect's last words use the same verb
and name the same deliverer: \latin{te liberánte
salvári}, ``that we may be saved, with thee setting us free.''

The Gradual's verse asks for help of another kind: ``make known to me thy
ways.'' St Robert Bellarmine reads that request as a prayer to be taught ``not
speculatively, but practically,'' that is, for ``such grace as may move his
will to observe the commandments cheerfully'' (\work{Explanatio}, Ps 24, on
v.~4).

William Durandus heard the waking in the chant. He notes that the cantor
begins \latin{Ad te levávi} low and lifts his voice, \latin{quod proprium est
excitantis}, ``which is the mark of one rousing others,'' and that the same is
marked in the Epistle, \latin{Hora est nos iam de somno surgere}
(\work{Rationale divinorum officiorum} VI, on the first Sunday of Advent).
Sicard of Cremona, writing of the Introit's ways and the Epistle, sets out
which are ways and which are wanderings: \latin{Via est de somno surgere,
arma lucis induere, honeste ambulare. Devium est in nocte pigritari, opera
tenebrarum operari}, ``The way is to rise from sleep, to put on the arms of
light, to walk becomingly. The wandering is to idle in the night and do the
works of darkness'' (\work{Mitrale} V, on the first Sunday of Advent).

The same antiphon returns at the offering: ``To thee have I lifted up my
soul.'' Aquinas reads the lifting as the preparation for prayer, which is, in
St John Damascene's words, \latin{ascensus intellectus in Deum}, ``the ascent
of the mind to God.'' The psalmist lifts his soul, Aquinas says, ``by
contemplation, looking on thy goodness, and by love,'' and he prays with
trust, as Paul says, ``through Christ, towards God. Not that we are
sufficient to think any thing of ourselves'' (2~Cor 3:4--5; \work{Super Psalmos} 24,
n.~1). Augustine describes his delay as a weight: \latin{sarcina saeculi,
velut somno adsolet, dulciter premebar}, ``I was sweetly weighed down by the
world's baggage, as one is by sleep'' (VIII.5.12). Sung after Paul's command,
the Offertory's lifted soul is the opposite of that weight, and its next
words, \latin{Deus meus, in te confido}, ``in thee, O my God, I put my
trust,'' are the trust that
Aquinas, with St Paul, places in God through Christ and not in ourselves.

\subsection{The day at hand and the kingdom near}

The Epistle says \latin{dies \dots\ appropinquávit}, ``the day is at hand,''
and the Gospel \latin{appropínquat redémptio vestra}, ``your redemption is at
hand,'' with the same verb. Sicard, writing of the churches that read this
Gospel, says that the Epistle is fitted (\latin{adaptatur}) to its ``the
kingdom of God is at hand'' by its ``our salvation is nearer,'' and Durandus
repeats him (\work{Mitrale} V; \work{Rationale} VI). The word Paul uses for
drunkenness, \latin{ebrietas}, returns in Christ's warning just after the
lesson against hearts weighed down \latin{in crapula, et ebrietate} (Lk
21:34). Christ and Paul both warn against a heart dulled by excess when the
day arrives.

Between the two lessons the Alleluia asks, \latin{salutáre tuum da nobis},
``grant us thy salvation'': the salvation Paul has just called nearer is
asked for as God's gift. Aquinas allows Paul's words one further sense. They
may be taken, he says, of \latin{tempus misericordiae}, ``the time of mercy,''
when a man begins to want to leave his past sins, for then he is nearer his
salvation than at the beginning, when his faith was not yet formed by charity
(\work{Super Epistolam ad Romanos} 13, lect.~3). If Paul's words bear that
sense, the Alleluia's first request, ``Shew us, O Lord, thy mercy,'' asks for
the moment Aquinas describes.

Paul makes the casting off of the works of darkness the hearers' own task.
The Secret asks God for a cleansing like it: that the
holy gifts, \latin{poténti virtúte}, ``by their mighty power,'' make us
\latin{mundátos} and \latin{purióres}, cleansed and purer.

At Communion the choir sings, ``The Lord will give goodness: and our earth
shall yield her fruit.'' In the same passage Sicard takes the goodness for a
good affection, the earth for our flesh and its fruit for good works, and
Durandus again follows him (\work{Mitrale} V; \work{Rationale} VI). The flesh
that Paul would keep from revelling and wantonness then bears the good works
by which, Aquinas said, a man rises from the sleep of negligence, and the
verse puts the Lord's gift before the earth's yield.

The Postcommunion looks to \latin{reparatiónis nostræ ventúra solémnia}, the
coming solemnities of our restoration, to be met \latin{cóngruis honóribus},
with fitting honours. The feast it looks to is that of the first coming, the
sense Aquinas gives Paul's words as the Church reads them in Advent.

\subsection{The difficulty: what kind of nearness}

Chrysostom's nearness is of time: each day brings resurrection and judgment
closer. Augustine, in his letter to Hesychius, refuses to measure that time
and holds that the desire for the Lord's coming is right whether he comes soon
or late (Letter 199.1). Aquinas's two senses, given above, let the Mass be
heard both ways: the Collect's \latin{veni} is addressed to the Son who came
at Bethlehem and will come again, and the Epistle's day is both the day of
glory Paul meant and the day of grace the season remembers.

\begin{fourSenses}
\item[Literal.] Paul tells the Romans that it is time to wake, because the
day of salvation is nearer than when they first believed, and that they are
to put away the works of darkness and put on the Lord Jesus Christ (Rom
13:11--14a).
\item[Allegorical.] As the Church reads the Apostle in Advent, the salvation
nearer than when men first believed is Christ in his first coming, and Paul
speaks for all the faithful since the world began (Aquinas); the night past
is the time before the Incarnation (Aquinas); the garment is Christ himself,
the King, who surrounds the soul and dwells in it (Chrysostom).
\item[Moral.] Rise from the sleep of sin by penance and from negligence by
good works (Aquinas); rouse yourself the more as the goal comes nearer, and
wake your neighbour (Chrysostom); and since the sleeper says ``presently,''
ask Christ who comes to set you free, as Augustine learned he could be freed
by grace alone (the Collect; Augustine's \work{Confessions}).
\item[Anagogical.] The day of glory comes nearer every day, and the one who
has put on Christ goes out to meet it as soldier and bridegroom (Chrysostom;
Aquinas on the Apostle's meaning).
\end{fourSenses}
